\section{余切丛}
\label{sec:cotangent-bundle}

我们首先注意到，哈密顿作用量泛函
\[
  \mathcal A_H(\lambda)
  =\int_\lambda p\,dq-\int_0^1 H\,dt
\]
可以推广到任意流形 $N$ 的\emph{余切丛} $T^*N$ 上。注意，对于逐纤维凸函数，勒让德变换仍可定义在切丛 $TN$ 上；因此，我们关于哈密顿作用量泛函的讨论也涵盖 $TN$ 上的拉格朗日作用量泛函。故下文将集中讨论余切丛。关于一般流形 $N$ 上拉格朗日动力系统的一份优美的数学论述，可参阅法蒂的著作 (Fa05)。

以 $\pi:T^*N\longrightarrow N$ 表示典范投影。设 $(q^1,\ldots,q^n)$ 是开集 $U\subset N$ 上任意一个局部坐标系；它在
$T^*U=T^*N|_U\subset T^*N$ 上诱导典范坐标
\[
  (q^1\circ\pi,\ldots,q^n\circ\pi,p_1,\ldots,p_n).
\]
只要不会引起混淆，我们就略微滥用记号，直接把 $q^j\circ\pi$ 仍记为 $q^j$。回忆，纤维坐标 $(p_1,\ldots,p_n)$ 是相对于 $T^*N$ 在开集 $U$ 上的局部标架
\[
  \{dq^1,\ldots,dq^n\}
\]
所取的线性坐标。考虑定义在 $T^*U$ 上的 1-形式
\begin{equation}
  \sum_j p_j\,d(q^j\circ\pi).
  \label{eq:2.1.1}
\end{equation}
下面这个引理的证明留作习题。

\begin{lemma}
\label{lem:liouville-form-global}
局部定义的 1-形式 \eqref{eq:2.1.1} 是整体定义的。更准确地说，设另一组典范坐标为
\[
  (Q^1\circ\pi,\ldots,Q^n\circ\pi,P_1,\ldots,P_n),
\]
其中 $(Q^1,\ldots,Q^n)$ 是 $V\subset N$ 上的局部坐标系，则在交叠 $U\cap V$ 上有
\[
  \sum_j p_j\,d(q^j\circ\pi)
  =\sum_j P_j\,d(Q^j\circ\pi).
\]
\end{lemma}

\begin{exercise}
证明引理 \ref{lem:liouville-form-global}。
\end{exercise}

这提示我们，应该可以对上述整体定义的 1-形式给出一个不依赖坐标的描述。

\begin{definition}[刘维尔 1-形式]
\stfhIndex{I197}{liu^wei^er^1^xing^shi}{刘维尔 1-形式（Liouville one form）}
\label{def:liouville-one-form}
在 $T^*N$ 上定义 1-形式 $\theta$ 如下：
\[
  \theta_\alpha(\xi)=\alpha(d\pi(\xi)),
\]
其中 $\alpha\in T^*N$，$\xi\in T_\alpha(T^*N)$。这个 1-形式称为\emph{刘维尔 1-形式}。
\end{definition}

立即可见，在典范坐标 $(q^1,\ldots,q^n,p_1,\ldots,p_n)$ 中，$\theta$ 恰好写成 \eqref{eq:2.1.1}。下列性质刻画了刘维尔 1-形式 $\theta$；其证明留作习题。

\begin{proposition}
\label{prop:liouville-characterization}
刘维尔形式 $\theta$ 满足下列性质。
\begin{enumerate}[label=(\arabic*)]
  \item 设 $\alpha$ 是 $N$ 上任意 1-形式，则
  \begin{equation}
    \widetilde\alpha^*\theta=\alpha
    \label{eq:2.1.2}
  \end{equation}
  在 $N$ 上成立。这里，等式左边的 $\widetilde\alpha:N\to T^*N$ 表示由丛 $\pi:T^*N\to N$ 的截面 $\alpha$ 所诱导的映射。
  \item 对任意垂直切向量 $\xi$，即满足 $d\pi(\xi)=0$ 的 $\xi$，都有 $\theta(\xi)=0$。
\end{enumerate}
反之，任何满足上述两个性质的 1-形式都等于 $\theta$。
\end{proposition}

\begin{proof}
由 $\theta$ 的定义很容易验证这两个性质。这里只证明逆命题。

设 $T^*N$ 上的 1-形式 $\theta'$ 满足上述两个性质。任取一点 $x\in T^*N$，并令 $q=\pi(x)$。选取 $N$ 上一个满足 $\widetilde\alpha(q)=x$ 的 1-形式 $\alpha$。于是
\[
  T_x(T^*N)=VT_x(T^*N)\oplus d\widetilde\alpha(T_qN),
\]
其中 $VT_x(T^*N)$ 是 $x\in T^*N$ 处的垂直切空间。性质 (2) 蕴含
\[
  \theta'(x)|_{VT_x(T^*N)}\equiv
  \theta(x)|_{VT_x(T^*N)}.
\]
因此，只需对所有 $v\in T_qN$ 证明
\[
  \theta'(x)(d\widetilde\alpha(v))
  =\theta(x)(d\widetilde\alpha(v)).
\]
而这个等式等价于
\[
  \widetilde\alpha^*\theta'(q)(v)
  =\widetilde\alpha^*\theta(q)(v).
\]
后一个等式显然成立，因为性质 (1) 给出
$\widetilde\alpha^*\theta'=\alpha=\widetilde\alpha^*\theta$。证明完毕。
\end{proof}

我们还要指出，性质 (1) 与 (2) 合起来蕴含
\begin{equation}
  \theta\equiv 0
  \label{eq:2.1.3}
\end{equation}
在零截面上成立。

事实表明，在更一般的情形中存在一种更基本的几何结构，它由 $\theta$ 的外微分诱导，称为辛形式。

\begin{definition}[典范辛形式]
\stfhSubIndex{I278}{xin^xing^shi}{辛形式（symplectic form）}{dian^fan^xin^xing^shi}{典范辛形式（canonical symplectic form）}
\label{def:canonical-symplectic-form}
设 $N$ 是任意流形。$T^*N$ 的\emph{典范辛形式}记为 $\omega_0$，定义为
\[
  \omega_0=-d\theta.
\]
\end{definition}

在 $d\theta$ 前选择负号，是为了使 $\omega_0$ 的坐标表达式为
\[
  \sum_j dq^j\wedge dp_j,
\]
而不是一些作者采用的 $\sum_jdp_j\wedge dq^j$。由定义，$\omega_0$ 是闭的（事实上是正合的），并在每个切空间 $T_\alpha(T^*N)$ 上给出一个斜对称双线性二形式。下面的非退化性则不那么显然。

\begin{lemma}
\label{lem:canonical-form-nondegenerate}
二形式 $\omega_0$ 是非退化的：在任意点 $\alpha\in T^*N$，对 $\xi\in T_\alpha(T^*N)$，有
\[
  \xi=0\quad\Longleftrightarrow\quad \xi\mathbin{\lrcorner}\omega_0=0.
\]
\end{lemma}

\begin{proof}
由于 $\omega_0$ 是张量，我们先在典范坐标 $(q^1,\ldots,q^n,p_1,\ldots,p_n)$ 中将它写出。由习题 2.1.2，
\[
  \omega_0=\sum_jdq^j\wedge dp_j,
\]
从中非退化性一目了然。
\end{proof}
